% Source: https://artofproblemsolving.com/community/c3h2615973_problem_4_of_finals__vi_ifym_sozopol_2015_67_theme
% Stable identity: aops-topic-2615973-q01
Write the number as $\overline{abcd}$ with $a\ne0$. Divisibility by $11$ means $(a+c)-(b+d)$ is a multiple of $11$. Also $(a+c)+(b+d)=22$, so the difference $(a+c)-(b+d)$ is even and lies between $-18$ and $18$. The only even multiple of $11$ in that range is $0$, so
\[a+c=b+d=11.\]
\begin{itemize}
\item Pairs $(a,c)$ of digits with $a+c=11$: $a=2,3,\dots,9$, so $8$ pairs (here $a\ne0$ automatically). Likewise $8$ pairs $(b,d)$. So $8\cdot8=64$ numbers satisfy the first two conditions.
\item Those without the digit $2$: a pair with sum $11$ contains a $2$ exactly when it is $(2,9)$ or $(9,2)$. So $6$ pairs for $(a,c)$ and $6$ for $(b,d)$ avoid the digit $2$, giving $36$ numbers.
\end{itemize}
Hence there are $64-36=\boxed{28}$ interesting numbers.
